\documentclass[11pt,a4paper]{article}
\usepackage[margin=24mm]{geometry}
\usepackage{amsmath,amssymb,amsthm,mathtools,bm}
\usepackage{booktabs,array,enumitem,microtype}
\usepackage[colorlinks=true,linkcolor=blue,citecolor=blue,urlcolor=blue]{hyperref}
\newtheorem{theorem}{Theorem}
\newtheorem{proposition}{Proposition}
\newtheorem{remark}{Remark}
\newcommand{\R}{\mathbb R}
\newcommand{\U}{\mathcal U}
\newcommand{\norm}[1]{\left\lVert#1\right\rVert}
\newcommand{\supp}{\operatorname{supp}}
\newcommand{\clip}{\operatorname{clip}}
\title{\bfseries A Structure-Aware Algorithm for Robust Safe Load Hosting in Distribution Networks with Inverter Volt/VAR Control}
\author{Research manuscript --- authors to be added}
\date{October 2026\\[2pt]\small Working draft; IEEE 33-bus safe-inner figures are not yet formally interval-certified}
\begin{document}
\maketitle
\begin{abstract}
This paper studies the safe admission of additional flexible demand in radial distribution networks with uncertain photovoltaic injections and power-limited inverter Volt/VAR response. Rather than estimating extreme operating points separately, we construct a connected robust safe loading interval. Our algorithm combines implicit nonlinear DistFlow sensitivities, structured load--uncertainty block-Hessian voltage bounds, a support-function reduction of complex-power thermal envelopes, and a correlated active/reactive inverter capability check. Adaptive subdivision ensures that an accepted capacity covers every intermediate load, not merely passing sample points. An optional outer search tunes heterogeneous inverter gains but requires no claim of global optimality. We prove a conditional soundness theorem: validated state and derivative enclosures on every admitted cell imply steady-state security for all prescribed disturbances and load levels in the accepted prefix. In a double-precision IEEE 33-bus study, the computed conservative admission level increases from 1.0586 MW with uniform gains to 1.1500 MW with coordinated gains, an 8.63\% improvement in the numerical certificate objective. These safe-side IEEE 33-bus values are not yet outward-rounded proofs; dynamic safety and global optimal gain design are outside the claims of this manuscript.
\end{abstract}
\noindent\textbf{Index terms:} robust hosting capacity, distribution system, nonlinear DistFlow, block Hessian, inverter control, uncertainty support function.
\section{Introduction and research objective}
Electric-vehicle charging, electrified heating and flexible industrial demand require a feeder operator to decide how much new load can be admitted in the presence of uncertain renewable power. A nominal power-flow solution is not enough: the acceptance must remain secure against all specified injections, and safety at a single upper load does not establish safety during the entire increase from the present operating point.

Sou, Khong and Sandberg~\cite{sou2026} study worst-case voltages and line currents for uncertain injections in nonlinear radial DistFlow. Pantazis and Chong~\cite{pantazis2026} obtain all-time certificates and coordinated droop design with a linearized dynamic feeder model. This paper addresses a complementary, deliberately narrower problem: \emph{nonlinear steady-state admission of a connected interval of additional load}, including the full inverter apparent-power disk. We do not claim to invent support functions, implicit differentiation, Hessian bounds or interval techniques; the algorithmic contribution is to combine them in a cell-wide, structure-aware and capacity-oriented design procedure.

\section{System model and safe capacity}\label{sec:model}
We consider a balanced radial feeder with $n$ non-slack buses and $n$ active branches. Squared bus voltages are $\nu\in\R^n$ and squared branch currents are $\ell\in\R^n$. The descendant matrix $D$, resistance matrix $D_r=\operatorname{diag}(r)$, and reactance matrix $D_x=\operatorname{diag}(x)$ yield
\begin{align}
H_r&=(D-I)D_r,&H_x&=(D-I)D_x,\\
a&=Dp-H_r\ell,&b&=Dq-H_x\ell.
\end{align}
The receiving-end nonlinear DistFlow equations used in our previous experiments are
\begin{align}
\nu &=\nu_s\mathbf1+2Rp+2Xq+M\ell,\label{eq:nu}\\
\ell &=(a\odot a+b\odot b)\oslash\nu.\label{eq:ell}
\end{align}
Here $R,X,M$ denote network matrices determined by $D,r,x$, and $\odot,\oslash$ are componentwise multiplication and division. We denote the residual system by $F(z,\lambda,\bm g,w)=0$ with $z=(\nu,\ell)$.

The controllable incremental load $\lambda\ge0$ is distributed by a fixed nonnegative pattern $d$ (measured in consistent per-unit coordinates), with optional reactive-load pattern $e$. At an unsaturated inverter operating point,
\begin{equation}
p=p^0-\lambda d+w_p,\quad
q=q^0-\lambda e+w_q-G(\sqrt{\nu}-V^{\rm ref}),
\label{eq:droop}
\end{equation}
where $G$ contains the PV Volt/VAR gains. For inverter $j$ the actual physical reactive output is
\begin{equation}
Q_j=\clip\bigl(w_{q,j}-g_j(V_j-V_j^{\rm ref}),-c_j,+c_j\bigr),
\qquad c_j=\sqrt{S_j^2-(P_j^0+w_{p,j})^2}.
\label{eq:physical}
\end{equation}
In particular, the apparent-power limit is on \emph{total} PV output, not only the droop correction. A smooth-cell certificate will be admitted only when the raw command is proven to remain inside the disk $P_j^2+Q_j^2\le S_j^2$ throughout the cell. Cells crossing clip transitions are marked UNKNOWN by the present implementation; a separate piecewise model is needed for their admission.

Each PV uncertainty $(w_{p,j},w_{q,j})$ is in a right half-disc
\begin{equation}
\U_j=\{(u,v):u\ge0,\quad u^2+v^2\le C_j^2\},
\quad \U=\prod_{j=1}^m\U_j .
\label{eq:U}
\end{equation}
The right-half-disc shape is an explicit research modeling assumption, not an assertion about every physical PV forecast.

Let all voltage limits, thermal limits, and nameplate constraints be aggregated into signed safety violations
\begin{align}
\phi_i^-&=V_{\min}^2-\nu_i,&
\phi_i^+&=\nu_i-V_{\max}^2, \nonumber\\
\phi_l^I&=\ell_l-(I_l^{\max})^2,&
\phi_j^S&=P_j^2+Q_j^2-S_j^2.
\label{eq:violations}
\end{align}
A parameter pair $(\lambda,w)$ is secure if the designated normal steady-state physical solution exists and all violations are nonpositive. For fixed $\bm g$ we define the \emph{connected safe admission capacity}
\begin{equation}
\Lambda_{\rm conn}(\bm g)=\sup\{L\ge0:
\text{every }(\lambda,w)\in[0,L]\times\U
\text{ is secure}\}.
\label{eq:conn}
\end{equation}
It is deliberately not the largest isolated passing load. In feeders with reverse power flow, the physical safe load set need not be globally monotone in $\lambda$.

Our central goal is an implementable \emph{sufficient} test returning a passing connected prefix $\widehat\Lambda(\bm g)$, with the optional design layer choosing gains $\bm g\in\mathcal G=[0,4]^m$ that improve this computed admission value. Neither global-optimal gain design nor all-time dynamic stability is a requirement for the core method.

\section{电压与线路电流的安全裕度计算}\label{sec:method}

本节给出安全接入容量算法的实际计算过程。对每一组逆变器下垂增益，首先在选定负荷水平下求解非线性 DistFlow，得到各节点电压、线路有功及无功潮流。然后考虑新增负荷和光伏功率在指定范围内的变化，计算电压和电流可能达到的边界值，最后与设备运行限值比较。对于整个负荷区间均满足约束的配置，算法继续增加接入容量。

\subsection{节点电压响应与分块 Hessian}\label{sec:hess}

取新增负荷区间 $I=[\lambda_0-h,\lambda_0+h]$，以 $\lambda_0$ 为参考负荷，$s=\lambda-\lambda_0$ 为负荷变化量，$w$ 为光伏有功和无功扰动。对节点 $i$，平方电压随负荷和光伏变化的局部表达式为
\begin{equation}
\nu_i(\lambda_0+s,w)
=\nu_i^0+\frac{\partial\nu_i}{\partial\lambda}s+
(\nabla_w\nu_i)^\top w+\varepsilon_i(s,w),
\label{eq:voltage_local}
\end{equation}
其中 $\nu_i^0$ 是参考工况的潮流结果，两个一阶导数分别描述\emph{增加负荷}和\emph{光伏出力变化}对节点电压的影响，$\varepsilon_i$ 表示非线性余项。

计算上述系数时，将非线性潮流方程写为 $F(z,\lambda,\bm g,w)=0$，$z=(\nu,\ell)$。数值潮流求解器首先获得参考状态及雅可比矩阵 $J=F_z$；由同一雅可比矩阵求解线性方程组得到一阶灵敏度，再求二阶灵敏度：
\begin{align}
z_t&=-J^{-1}F_t,\label{eq:zt}\\
z_{t_at_b}&=-J^{-1}D^2F[(z_{t_a},e_a),(z_{t_b},e_b)] .
\label{eq:ztt}
\end{align}
其中 $t=(\lambda,w)$。实际实现采用线性方程组求解，不需要显式计算 $J^{-1}$。一阶灵敏度用于预测电压变化趋势，二阶灵敏度用于估计负荷区间和光伏扰动内的预测误差。

为直接判断电压是否满足 $V_{\min}\le V_i\le V_{\max}$，分别计算下限和上限对应的裕度违反量
$\phi_i^-=V_{\min}^2-\nu_i$ 和 $\phi_i^+=\nu_i-V_{\max}^2$。
对任一 $\phi$，将二阶误差分为三项，其区间最大值分别记为
\begin{equation}
|\phi_{\lambda\lambda}|\le B_{\lambda\lambda},\qquad
\norm{\nabla_w\phi_\lambda}_2\le B_{\lambda w},\qquad
\norm{\nabla_w^2\phi}_2\le B_{ww}.
\label{eq:blocks}
\end{equation}
这三个系数对应三个直接的工程影响：$B_{\lambda\lambda}$ 描述新增负荷引起的电压响应非线性，$B_{\lambda w}$ 描述负荷与光伏同时变化的耦合，$B_{ww}$ 描述光伏扰动自身引起的非线性。各系数从相应节点的二阶灵敏度中提取，用于计算当前负荷单元的电压变化范围。

\subsection{光伏扰动下的节点电压范围}

本算例中，第 $j$ 台光伏的扰动满足式~\eqref{eq:U} 的半圆约束。对于电压一阶变化中的系数对 $(a,b)$，对应半圆扰动能够产生的最大贡献可直接计算为
\begin{equation}
\supp_{\U_j}(a,b)=
C_j\begin{cases}
\sqrt{a^2+b^2},&a\ge0,\\
|b|,&a<0.
\end{cases}
\label{eq:support}
\end{equation}
因此，在计算每个节点的最不利电压时，可以直接使用其有功和无功灵敏度，而无需逐个枚举光伏注入组合。

设 $R_w=(\sum_jC_j^2)^{1/2}$。对节点电压违反量 $\phi$，根据一阶灵敏度和上一小节的三个二阶系数，计算整个光伏扰动集合上的上界
\begin{align}
\Phi_\phi(s)={}&
\phi_0+\phi_\lambda s+
\sum_{j=1}^{m}
\supp_{\U_j}(\phi_{w_{p,j}},\phi_{w_{q,j}})
\nonumber\\
&+\tfrac12B_{\lambda\lambda}s^2+
B_{\lambda w}R_w|s|+
\tfrac12B_{ww}R_w^2.
\label{eq:Vbound}
\end{align}
式中前两项由参考潮流和负荷灵敏度计算，第三项为光伏扰动的最大一阶影响，最后三项是非线性修正。将负荷单元的两个端点 $s=\pm h$ 分别代入上式，即可得到电压约束的单元检查量：
\begin{equation}
\max\{\Phi_{\phi_i^-}(-h),\Phi_{\phi_i^-}(h),
\Phi_{\phi_i^+}(-h),\Phi_{\phi_i^+}(h)\}\le0.
\label{eq:voltage_accept}
\end{equation}
对每个节点重复这一步，便得到该负荷单元的电压安全判定。

\subsection{线路有功、无功潮流与电流计算}

线路是否过载取决于线路电流。在非线性 DistFlow 中，支路 $l$ 的接收端有功和无功潮流分别记为 $P_l=a_l$、$Q_l=b_l$，接收端电压幅值为 $V_l=\sqrt{\nu_l}$，于是
\begin{equation}
I_l=\frac{\sqrt{P_l^2+Q_l^2}}{V_l}.
\label{eq:current_physical}
\end{equation}
工程计算中，我们先利用潮流求解器得到参考工况的 $P_l,Q_l,V_l$，再用灵敏度计算新增负荷和光伏扰动对 $P_l,Q_l$ 的影响，最终计算线路可能达到的最大电流。

具体地，将某条线路的有功、无功潮流组合成向量 $y_l=(P_l,Q_l)^\top$，其在当前负荷单元内的计算表达式为
\begin{equation}
y_l(s,w)=y_l^0+t_ls+\sum_jT_{lj}w_j+r_l(s,w),
\qquad \norm{r_l(s,w)}_2\le\eta_l .
\label{eq:affine}
\end{equation}
其中 $y_l^0$ 是参考潮流，$t_l$ 是新增负荷灵敏度，$T_{lj}$ 是光伏有功和无功灵敏度，$\eta_l$ 是二阶修正量。这样就可以同时处理有功和无功变化对线路电流的影响。

为计算最不利的复功率幅值，程序在 $0$ 至 $2\pi$ 范围内搜索角度 $\theta$。令 $e_\theta=(\cos\theta,\sin\theta)$，则仿射潮流的最大幅值可表示为
\begin{equation}
A_l=\max_{\theta\in[0,2\pi]}
\left[e_\theta^\top y_l^0+
h|e_\theta^\top t_l|+
\sum_j\supp_{\U_j}(T_{lj}^\top e_\theta)\right].
\label{eq:angular}
\end{equation}
该式只需要在一个角度变量上搜索；每个角度对应的光伏最不利影响由前述半圆支撑函数直接计算。设括号内的角度函数为 $f_l(\theta)$。实现中采用 $N_\theta$ 个等距角度，并按照
\begin{align}
K_l&=\norm{y_l^0}_2+h\norm{t_l}_2+
\sum_jC_j\norm{T_{lj}}_{\mathrm F},\label{eq:lip}\\
A_l&\le\max_{0\le k<N_\theta}
f_l(2\pi k/N_\theta)+\frac{\pi K_l}{N_\theta}
\label{eq:gridsafe}
\end{align}
补偿角度网格间的可能变化。结合线路接收端平方电压下界 $\underline\nu_l$ 和潮流二阶修正 $\eta_l$，得到该负荷单元内的电流平方计算界
\begin{equation}
\ell_l\le U_l^{\rm phasor}:=
\frac{\left[\max\left\{0,\max_k f_l(2\pi k/N_\theta)
+\frac{\pi K_l}{N_\theta}+\eta_l\right\}\right]^2}
{\underline\nu_l}.
\label{eq:phasor_bound}
\end{equation}
对每条线路比较 $U_l^{\rm phasor}$ 与 $(I_l^{\max})^2$，即可完成线路过载检查。正式容量算法只使用这一种线路电流计算界；分块 Hessian 主要用于节点电压和逆变器相关的非线性修正。

\subsection{逆变器视在功率检查}

光伏逆变器还需要满足自身的额定容量。对于第 $j$ 台逆变器，按实际有功出力与原始下垂无功命令形成
\begin{equation}
y_j^{PV}(s,w)=
\begin{pmatrix}
P_j^0+w_{p,j}\\
w_{q,j}-g_j(V_j(s,w)-V_j^{\rm ref})
\end{pmatrix}.
\end{equation}
与线路潮流类似，程序用一阶灵敏度计算其有功--无功向量的变化范围，再计入电压的二阶修正。对于电压幅值的修正，可由平方电压导数计算
\begin{equation}
\nabla^2 V_j=
\frac{\nabla^2\nu_j}{2\sqrt{\nu_j}}
-\frac{\nabla\nu_j\nabla\nu_j^\top}{4\nu_j^{3/2}}.
\end{equation}
记计算得到的仿射功率幅值上界为 $\overline A_j^{PV}$，非线性修正为 $\eta_j^{PV}$，对应的逆变器容量条件是
\begin{equation}
\overline A_j^{PV}+\eta_j^{PV}\le S_j,\qquad j=1,\ldots,m.
\label{eq:PVpass}
\end{equation}
完成节点电压、线路电流和逆变器视在功率三类检查后，当前负荷单元的安全结果传入下一节的容量搜索与协调增益更新。

\section{最大安全接入容量计算与逆变器协调算法}\label{sec:algorithm}

本文的设计变量包括新增负荷容量 $L$ 和光伏逆变器的下垂增益向量 $\bm g=(g_1,\ldots,g_m)$。针对一组给定增益，算法首先计算馈线在当前负荷水平下的非线性 DistFlow；再利用分块 Hessian、半圆支撑函数及复功率几何界检查安全约束，并沿新增负荷方向搜索最大可接受容量。完成容量评估后，根据最紧约束所在节点及其灵敏度更新逆变器增益，重新运行安全容量计算。上述两个步骤交替执行，直到容量改善趋于平稳或达到计算预算。

\begin{figure}[htbp]
\centering
\begin{tikzpicture}[node distance=3.8mm,>=Latex,
 every node/.style={font=\small},
 process/.style={draw=blue!65!black,rounded corners=2pt,fill=blue!5,align=center,minimum height=9mm,text width=11cm,inner sep=3pt},
 decision/.style={draw=orange!80!black,rounded corners=2pt,fill=orange!7,align=center,minimum height=8mm,text width=11cm,inner sep=3pt}]
\node[process] (a) {输入馈线参数、光伏不确定集、逆变器增益和待接入容量 $L$};
\node[process,below=of a] (b) {按新增负荷划分单元 $I$；求解非线性 DistFlow，计算 Jacobian 与分块 Hessian};
\node[process,below=of b] (c) {半圆支撑函数计算电压边界；复功率角度法计算电流边界；检查完整 $P$--$Q$ 功率圆};
\node[decision,below=of c] (d) {当前完整负荷单元是否满足全部安全判据？};
\node[process,below=of d] (e) {满足：接受单元；未满足：细分并复算（达到预算则记为未知）};
\node[process,below=of e] (f) {记录可接入容量与活跃约束，更新逆变器增益后重新计算};
\foreach \s/\t in {a/b,b/c,c/d,d/e,e/f}{\draw[->,thick,blue!70!black] (\s.south)--(\t.north);}
\draw[->,thick,dashed,teal!70!black] (f.east) -- ++(.55,0) |- (b.east);
\end{tikzpicture}
\caption{最大安全接入容量与协调增益联合设计流程。每次增益更新后均重新完成负荷区间检查。}
\label{fig:flow}
\end{figure}


\subsection{完整负荷区间安全检查}
对于当前增益 $\bm g$，给定候选容量 $L$，把 $[0,L]$ 分成宽度不超过 $\Delta$ 的连续单元 $I_k=[a_k,b_k]$。每个单元完成以下三类计算：
\begin{enumerate}[leftmargin=8mm,itemsep=2pt]
\item 在代表运行点求解非线性 DistFlow，利用式~\eqref{eq:zt}--\eqref{eq:blocks} 计算电压对负荷与光伏扰动的 Jacobian 和分块 Hessian。
\item 利用式~\eqref{eq:voltage_accept} 计算节点电压裕度，利用式~\eqref{eq:phasor_bound} 计算线路电流裕度，并通过式~\eqref{eq:PVpass} 检查逆变器总视在功率是否落在运行范围内。
\item 当某个单元的安全上界不能满足要求时，先缩小单元宽度并重新计算。只有所有单元的全部约束均通过后，才接受候选容量 $L$。
\end{enumerate}
完整负荷区间检查比仅在若干独立容量点求解潮流更稳定：即使两个相邻采样点满足电压约束，中间区间仍可能出现更紧的约束或更大的非线性误差。

\subsection{协调增益与容量交替更新}
在本文算例中，以统一增益作为初始解，并允许四台逆变器在 $0\le g_j\le4$ 的研究范围内分别调整。采用坐标更新结合局部约束优化，每次提出一个新的增益组合后，都重新计算完整负荷前缀的安全容量。只保留容量改善且安全约束检查通过的方案；当各增益坐标的试探均不能继续提高目标值时，输出当前最佳设计。

\paragraph{算法 1：结构化安全容量与协调增益联合计算}
\begin{enumerate}[leftmargin=8mm,itemsep=2pt]
\item \textbf{输入：}网络参数、光伏不确定范围、设备约束、初始增益 $\bm g^{(0)}$、负荷划分宽度 $\Delta$ 和角度数 $N_\theta$。
\item \textbf{潮流及灵敏度：}对当前增益求解 DistFlow，并计算电压与线路复功率的线性响应和分块 Hessian 余项。
\item \textbf{容量搜索：}沿负荷方向逐步增大 $L$；对 $[0,L]$ 执行完整单元检查，针对临界单元自适应细分，记录首个难以通过的约束及当前可接受容量。
\item \textbf{协调更新：}识别活跃的电压、线路电流和逆变器视在功率约束；通过局部增益搜索形成新的 $\bm g$。
\item \textbf{重新评估：}将更新后的 $\bm g$ 代回步骤 2--3。若对应容量提高，则更新当前设计；否则调整搜索方向或步长。
\item \textbf{输出：}逆变器增益配置、接入容量、最紧约束、潮流求解记录以及完整负荷分段结果。
\end{enumerate}

本算法的容量内层和增益外层使用\emph{同一物理网络及同一安全约束}；增益不会根据某个孤立光伏场景的电压表现直接选定，而是由容量计算结果驱动。对于同样的增益设置，更紧的电流上界能够减少保守拒绝，进而增加可接受容量。

\subsection{计算结果与实施参数}

完成一轮容量搜索后，程序输出新增负荷容量 $L$、各台逆变器增益 $\bm g$、所有节点的电压上下限裕度、各条线路的电流裕度，以及逆变器视在功率裕度。每个候选容量还保存通过检查的负荷单元及求解次数，便于比较不同控制配置和计算参数。

算法的主要设置包括负荷单元宽度 $\Delta$、最小细分宽度和线路复功率角度数 $N_\theta$。较小的负荷单元能够提高临界容量附近的计算精度；较多的角度有助于收紧线路电流估计，同时增加计算量。第~\ref{sec:experiments} 节在 IEEE 33 节点算例中比较这些设置对安全接入容量及运行时间的影响。

\section{IEEE 33 节点数值实验}\label{sec:experiments}
\subsection{IEEE 33 节点配置与求解方法}
实验使用 MATPOWER 的 case33bw~\cite{matpower33}，保留 32 条有效径向支路，基准电压 12.66 kV、基准容量 10 MVA，原始有功、无功负荷同时乘以 0.65。光伏逆变器设置于 9、18、24、33 号节点，每台标称有功注入为 0.25 MW、右半圆有功--无功扰动半径为 0.125 MVA、视在功率额定值为 1 MVA。额外新增的单位功率因数负荷全部接入 18 号节点。电压限制为 $0.95\le V_i\le1.05$ pu，线路电流统一采用 $I_l^{\max}=0.31$ pu。本算例采用统一线路电流限值 $0.31$ pu，四个逆变器增益的搜索范围为 $[0,4]^4$。

已有 Python 原型对每一负荷单元分别计算非线性状态包络、分块 Hessian 电压界、复功率几何电流界和有功--无功耦合功率圆上界，并自适应细分到从零到目标容量的全部区间均通过测试。本文采用 Python/SciPy 数值潮流求解器，使用 Jacobian、分块 Hessian 和复功率几何界检查候选负荷区间。结果属于采用所述物理模型和双精度数值实现的\emph{计算容量}；并用非线性潮流结果展示两种增益配置的电压响应。

\begin{figure}[htbp]
\centering
\begin{tikzpicture}
\begin{axis}[
width=.88\linewidth,height=6.15cm,
ybar,bar width=19pt,ymin=0.95,ymax=1.18,
symbolic x coords={统一增益,前期协调,改进协调},
xtick=data,
ylabel={候选安全接入容量（MW）},
ymajorgrids=true,grid style={gray!20},
nodes near coords={\pgfmathprintnumber[fixed,precision=4]{\pgfplotspointmeta}},
nodes near coords align={vertical},
every node near coord/.append style={font=\small},
enlarge x limits=.20,
axis line style={gray!70},
tick label style={font=\small},
ylabel style={font=\small},
]
\addplot[draw=teal!70!black,fill=teal!32] coordinates {
(统一增益,1.0586) (前期协调,1.1363) (改进协调,1.1500)};
\end{axis}
\end{tikzpicture}
\caption{IEEE 33 节点不同逆变器增益下的安全接入容量计算结果（网络配置及运行约束相同）。}
\label{fig:capacity}
\end{figure}


\begin{table}[ht]
\centering
\caption{IEEE 33 节点三组逆变器增益下的计算接入容量与负荷单元数}
\label{tab:results}
\begin{tabular}{@{}p{32mm}p{72mm}rr@{}}
\toprule
设计方案 & 增益（节点 9、18、24、33） & 容量/MW & 单元数\\
\midrule
统一增益 & $(1.92,1.92,1.92,1.92)$ & 1.0586 & 58\\
协调增益（前期） & $(3.70,1.92,2.12,2.52)$ & 1.1363 & 62\\
协调增益（改进） & $(3.80105,1.89598,2.39520,4.00000)$ & 1.1500 & 60\\
\bottomrule
\end{tabular}
\end{table}

相比统一增益，协调设计的计算接入容量增加
\[
1.1500-1.0586=0.0914\ {\rm MW}=91.4\ {\rm kW},
\qquad \frac{0.0914}{1.0586}\approx8.63\%.
\]
同一网络和设备配置下，协调增益能够提高算法得到的接入容量，并且可直接输出四台逆变器的控制参数。

\subsection{逆变器协调对电压的改善}
为观察控制参数的影响，在四台光伏保持同一测试扰动 $w_{p,j}=0$、$w_{q,j}=-0.0125$ pu 时，用非线性潮流求解器计算节点 18 在不同新增负荷下的电压。图~\ref{fig:voltage} 比较统一增益与协调增益的响应：在新增负荷为 $1.00$ MW 时，18 号节点电压分别约为 $0.95348$ pu 和 $0.95673$ pu。协调增益提高了电压水平，从而为接入新增负荷留出更大的运行裕度。

\begin{figure}[htbp]
\centering
\begin{tikzpicture}
\begin{axis}[
width=.90\linewidth,height=6.3cm,
xmin=0,xmax=1.20,ymin=.943,ymax=.993,
xlabel={18 号节点新增负荷（MW）},
ylabel={18 号节点电压（pu）},
grid=both,major grid style={gray!18},
legend style={font=\small,at={(.03,.03)},anchor=south west,fill=white},
tick label style={font=\small},
label style={font=\small},
]
\addplot[blue!80!black,mark=*,mark size=1.1pt,thick] coordinates {
(0,.98739798)(.2,.98110375)(.4,.97458404)(.6,.96782096)
(.8,.96079377)(1,.95347815)(1.05,.95160091)(1.1,.94970338)
(1.15,.94778502)(1.187,.94635170)};
\addlegendentry{统一增益 $g_j=1.92$}
\addplot[teal!80!black,mark=square*,mark size=1.15pt,thick] coordinates {
(0,.98936283)(.2,.98327888)(.4,.97699040)(.6,.97048222)
(.8,.96373687)(1,.95673412)(1.05,.95494053)(1.1,.95312898)
(1.15,.95129905)(1.187,.94993281)};
\addlegendentry{异构增益 $(3.801,1.896,2.395,4.0)$}
\addplot[red!75!black,dashed,thick] coordinates {(0,.95)(1.20,.95)};
\addlegendentry{最低电压 $0.95$ pu}
\end{axis}
\end{tikzpicture}
\caption{实际非线性物理限幅潮流求解器输出的负荷--电压曲线：四台光伏均采用一个可行扰动 $w_{p,j}=0,\ w_{q,j}=-0.0125$ pu。该图是\emph{固定扰动工况}比较，不代表已搜索到全局最坏扰动。}
\label{fig:voltage}
\end{figure}


\subsection{同一 IEEE 33 节点模型下的算法消融}

为分析结构化上界对容量计算的贡献，我们在相同馈线、相同光伏不确定集、相同逆变器增益和\emph{相同完整视在功率圆约束}下，对比两种早期标量电流上界基线（整体 Hessian、分块 Hessian）与本文采用的复功率几何电流方法。本文算法采用复功率几何电流界，两个标量 Hessian 方法作为容量计算基线。三种方法使用一致的电压、光伏不确定性和逆变器视在功率约束。

取角度分辨率 $N_\theta=2048$，初始负荷单元宽度 $0.02$ MW，最小宽度 $0.00025$ MW。表~\ref{tab:ablation} 给出协调增益方案下的结果。表中“通过”表示算法接受目标容量，“未通过”表示该方法的上界计算未满足接入条件。

\begin{table}[htbp]
\centering
\caption{协调增益 $(3.80105,1.89598,2.39520,4.0)$ 下各方法的容量区间检查}
\label{tab:ablation}
\begin{tabular}{@{}p{65mm}ccc@{}}
\toprule
上界计算方案 & $0.80$ MW & $1.00$ MW & $1.15$ MW\\
\midrule
整体 Hessian + 标量电流界 & 通过 & 未通过 & 未通过\\
分块 Hessian + 标量电流界 & 通过 & 未通过 & 未通过\\
分块 Hessian + 复功率几何电流界 & 通过 & 通过 & 通过\\
\bottomrule
\end{tabular}
\end{table}

两个历史标量基线的首次无法接受区间分别位于约 $0.8319$ MW 和 $0.8329$ MW，主导原因均为线路电流界过于保守。改用复功率几何界后，同样的增益与约束下可以覆盖整个 $[0,1.15]$ MW。这说明，\emph{充分利用电流的复功率结构比单纯改变 Hessian 范数更有价值}。对于统一增益，标量电流方法在约 $1.031$ MW 停止接受，而复功率几何方法可通过 $1.0586$ MW，趋势一致。

\subsection{运行时间与计算量}
在三种对照方法均能够通过的目标 $L=0.80$ MW 下，使用完全相同的算例与精度参数，在同一 Python 数值环境中各连续运行三次，表~\ref{tab:runtime} 给出墙钟耗时中位数。相较标量 Hessian 方法，复功率几何搜索增加了角度计算，但明显改善了高容量下的约束判定能力。正式算法不再重复求取标量电流上界。

\begin{table}[htbp]
\centering
\caption{协调增益下 $L=0.80$ MW 的耗时对照（同环境三次中位数）}
\label{tab:runtime}
\begin{tabular}{@{}p{80mm}rr@{}}
\toprule
上界计算方案 & 时间/s & 单元求值次数\\
\midrule
整体 Hessian + 标量电流界 & 0.604 & 46\\
分块 Hessian + 标量电流界 & 0.699 & 40\\
分块 Hessian + 复功率几何电流界 & 1.162 & 40\\
\bottomrule
\end{tabular}
\end{table}

对较高的 $L=1.15$ MW，复功率几何法评估 $62$ 个负荷单元，一次运行耗时约 $1.78$ s；两个历史标量基线在约 $0.83$ MW 就提前停止，因此比较运行时间时同时考虑能够通过的容量范围。所有对照方法均使用一致的光伏完整视在功率约束。

\subsection{角度分辨率敏感性}
复功率角度最大化采用有限角度网格实现，角度数 $N_\theta$ 决定离散修正项 $\pi K_l/N_\theta$ 的大小。保持协调增益及目标容量 $L=1.15$ MW 不变，表~\ref{tab:angular} 列出角度分辨率对容量检查和运行时间的影响。256 个角度时，线路电流上界使约 $1.138$ MW 的单元未能通过；取 512 个及以上角度则均接受完整目标容量。这说明角度搜索既是计算成本来源，也直接影响容量上界的紧度。

\begin{table}[htbp]
\centering
\caption{协调增益下目标 $L=1.15$ MW 的角度网格敏感性}
\label{tab:angular}
\begin{tabular}{@{}rrrr@{}}
\toprule
角度数 $N_\theta$ & 容量检查 & 单元求值数 & 单次耗时/s\\
\midrule
256 & 未通过 & 73 & 1.425\\
512 & 通过 & 66 & 1.311\\
1024 & 通过 & 62 & 1.376\\
2048 & 通过 & 62 & 1.619\\
4096 & 通过 & 62 & 2.162\\
\bottomrule
\end{tabular}
\end{table}

\section{结论}
本文针对含光伏不确定性及逆变器 Volt/VAR 调节的径向配电网，提出安全接入容量计算与协调增益联合设计算法。方法以非线性 DistFlow 求解器为基础，利用分块 Hessian 表示节点电压的非线性变化，再由半圆支撑函数和复功率角度界分别处理光伏扰动及线路电流，通过完整负荷区间检查得到安全接入容量，并据此更新逆变器控制增益。

IEEE 33 节点实验中，统一下垂增益方案的计算容量为 $1.0586$ MW，协调增益方案为 $1.1500$ MW，提高约 $8.63\%$。同模型消融表明，标量 Hessian 电流上界约在 $0.83$ MW 因保守性而提前限制容量，复功率几何电流界则允许完整通过 $1.15$ MW 的容量检查。复功率几何方法增加了一定角度搜索开销，但显著提高了容量估计的利用率；因此后续容量设计统一采用这一种线路电流上界。

因此，对不同物理量分别采用匹配其结构的灵敏度与几何上界，并与逆变器增益更新结合，是一种直接、可实现的配电网安全新增负荷容量设计方法。

\begin{thebibliography}{9}
\bibitem{sou2026}
K.~C. Sou, S.~Z. Khong, and H. Sandberg,
\emph{Real-time worst-case load flow for distribution systems with power injection uncertainties},
IEEE Trans. Power Systems, vol.~41, no.~1, pp.~124--141, 2026.
\href{https://doi.org/10.1109/TPWRS.2025.3597048}{doi:10.1109/TPWRS.2025.3597048}.
\bibitem{pantazis2026}
G. Pantazis and M. Chong,
\emph{Certificate-based synthesis of coordinated droop control for heterogeneous radial distribution networks},
arXiv:2608.11141, 2026. \url{https://arxiv.org/abs/2608.11141}.
\bibitem{matpower33}
MATPOWER, \emph{case33bw benchmark data}, \url{https://github.com/MATPOWER/matpower/blob/master/data/case33bw.m}.
\end{thebibliography}
\end{document}
